% !TeX root = ../main.tex
\chapter{INCOGNITOS: A Practical Unikernel Design for Full-System Obfuscation in Confidential Virtual Machines}\label{chap:incognitos}

\emph{* The content of this chapter were published in 2025 IEEE Symposium on Security and Privacy (SP), San Francisco, CA, USA, 4192-4209.}
\\



\epigraph{
	... abstraction are not designed to keep secrets.
}{\textit{Butler Lampson}~\cite{lampson2015perspectives}}

\epigraph{
	If you know the enemy and know yourself, you need not fear the result of a hundred battles.
}{\textit{Sun Tzu}, The Art of War}

\

\

\clearpage


To regain customers' trust in cloud, cloud providers are increasingly adopting specialized hardware from compute vendors (e.g., Intel, AMD).
These are new processors with secure primitives built into the silicon that can create a \gls{tee}, which protects sensitive data and programs from the rest of the platform, including compromised cloud administrators.

\gls{cvm} is the the next development in \gls{tee} technologies.
As cloud computing increasingly adopts  for trusted execution, the system’s traditional trust hierarchy is being inverted.
Where the operating system once resided within the trusted computing base (\gls{tcb}), it now operates alongside or even below untrusted hypervisors and shared hardware resources.

Consequently, even with hardware-provided isolation, side-channel attacks continue to undermine confidentiality guarantees.
This chapter introduces \incognitos, a unikernel-based operating system that restores coherence between trust and protection by reimagining the OS as a self-obfuscating execution environment within a CVM.
By retrofitting scheduling and memory management with adaptive obfuscation mechanisms, \incognitos transforms the OS from a potential attack surface into an active contributor to system-level confidentiality.
The remainder of this chapter presents the motivation, design, and evaluation of \incognitos, demonstrating how trust can be reestablished through self-obfuscation and dynamic adaptation in hardware-enforced environments.

\clearpage


\section{Toward coherent OS-assisted CVM protection}


\gls{cvm} brings unprecedented compatibility for cloud workloads, allowing users make an existing workload confidential.
This improves usability, reduces the adoption costs, enabling trivial protection of existing applications without modification, and cutting-edge applications like AI inference~\cite{yifan2025pipellm}.
Unfortunately, the \gls{cvm} also model inherit existing side-channel attacks challenges from the previous model of intel SGX, while also introducing unique challenges, due to the expanded \gls{tcb}.

\hdr{Side-channel attacks}
In 1973, Butler Lampson~\cite{lampson1973confinement} described in his seminar work, the difficulties of preventing an untrusted program from leaking secrets through \emph{covert channels} or \emph{side channels}.
These channels, such as the time an operation takes, while not designed for information transfer, allow a program to communicate with other programs.


In modern computing systems, side channels are major challenges in ensuring the security of \glspl{tee} and \gls{cvm}.
Previous research has demonstrate feasibilty of microarchitectural attacks such as cache attacks~\cite{yarom_flushreload_nodate,liu_last-level_2015},
that can be used to extract sensitive information.
\gls{tee} deployments excabates the problem with a much stronger threat model compare to traditional protections (\cref{sec:intro:bg-os}).
The attackers, now given access to powerful capabilities like page access interception~\cite{xu2015controlledchannel} and interrupt timers ~\cite{sgx-step,bulck2018nemesis,cachezoom}, are able extract sensitive from much more effectively.
Side channels were a significant problem with Intel SGX, and continue to persist in the \gls{cvm} model.
CVMs remain vulnerable to a broad class of side-channel attacks, including page-fault channels, cache-timing channels.
Worse, newer \gls{cvm}-specific side-channels, such as the ciphertext side channel~\cite{li_cipherleaks_nodate,ciphersteal}, continued to be discovered.


% On the other hand, microarchitectural side-channel attacks leverages the subtle difference in the microarchitectural state that are produced when processing secret information.


% Another challenge arises from the narrow interface between the guest and hypervisor; misunderstanding or misusing this interface can introduce vulnerabilities, as prior work has demonstrated through interface-level attacks.
% Finally, the OS remains responsible for resource management inside the VM yet lacks visibility into how specific operations strengthen or weaken enclave-level security.

% \hdr{Trusted I/O}
%
% The protection boundary established by the CPU does not extend to accelerators or I/O devices, which remain outside the TEE.
% 	[[PLACEHOLDER: describe why this is a problem or give an example attack scenario]]

\hdr{Taming the bloated TCB}

CVMs have enabled organizations to migrate large, modern applications—such as LLM inference workloads—into TEEs with minimal performance overhead and engineering cost.
However, doing so places a significant amount of code inside the TEE, and the security of the user's data now depends on large, feature-rich OSes such as Linux or Windows.
These OSes contain tens of millions of lines of code and thousands of known vulnerabilities.
Thus, while placing full OSes in the TEE simplifies deployment, it also increases the risk of compromise.

% Toward reducing the TCB, prior work has explored several approaches.
% Panoply~\cite{shinde_panoply_2017} decomposes enclave applications into smaller least-privilege components.
% Gramine-TDX~\cite{dimitrii2024graminetdx}, and serverless enclave systems~\cite{jiacheng2025serverless} offlload non-essential functionalities from the guest into the host to reduce the code size.




\hdr{\incognitos: coherent OS-asssited CVM protection}

\incognitos is designed to tackles the above challenges.
It explores the principle of coherent protection (\cref{chap:intro}) by rethinking the traditional protection model of \gls{os} (\cref{sec:intro:bg-os}), to align with the protection estalished by \gls{cvm} hardware like AMD SEV and Intel TDX.

This shift in perspective has two implications.
First, this means that the \gls{os} must now incorporate side-channel attacks from untrusted hypervisors into its threat model.
with the \gls{os}'s access to low-level capabilities like memory management and scheduling, is able to \emph{efficiently} and \emph{transparently} protect the memory accesses from side-channel attack.
Second,  it motivates the exploration of alternative OS designs to achieve the security goals.
With a unikernel design, \thename render whole-system protection feasible, while also keeping the TCB minimal.

\clearpage


% \newcommand*{\thename}{\textsc{IncognitOS}\xspace}
% Oblivious Unikernel
% Unikfraft Oblivious Page Oblivious SEV CVM
% ObliviUN 
% \newcommand*{\obfusmem}{\textsc{ObfusMem}\xspace}
% \newcommand*{\obfusmemch}{\textsc{ObfusMem-CH}\xspace}
% \newcommand*{\obfuscipher}{\textsc{ObfusCipher}\xspace}


\newacronym{mpk}{MPK}{Memory Protection Key}
\newacronym{idt}{IDT}{Interrupt Descriptor Table}
\newacronym{ovas}{OVAS}{ORAM-backed Virtual Address Space}
\newacronym{cfg}{CFG}{Control Flow Graph}
\newacronym{oee}{OEE}{OEE}
% \newacronym{ir}{IR}{Intermediate Representation}
% \newacronym{mmu}{MMU}{Memory Management Unit}
% \newacronym{mir}{MIR}{Machine IR}
% \newacronym{tee}{TEE}{Trusted Execution Environment}
% \newacronym{csp}{CSP}{Cloud Service Provider}
% \newacronym{sev}{SEV}{Secure Encrypted Virtualization}
% \newacronym{tdx}{TDX}{Trust Domain Extensions}
% \newacronym{cca}{CCA}{Confidential Compute Architecture}
% \newacronym{sgx}{SGX}{Software Guard Extension}
% \newacronym{vm}{VM}{Virtual Machine}
% \newacronym{ghcb}{GHCB}{Guest-Host Control Block}
\newacronym{msr}{MSR}{Model-Specific Register}
% \newacronym{cvm}{CVM}{Confidential VM}
\newacronym[plural=PTEs,firstplural=Page Table Entries]{pte}{PTE}{Page Table Entry}
% \newacronym{tlb}{TLB}{Translation-Lookaside Buffer}
\newacronym{pfn}{PFN}{Page Frame Number}
\newacronym{tsc}{TSC}{Time Stamp Counter}
% \newacronym{npf}{NPF}{Nested Page Fault}
\newacronym{npfongpt}{NPF-on-GPT}{Nested Page Fault on GPT}
% \newacronym{gpt}{GPT}{Guest Page Table}
% \newacronym{pf}{PF}{Page Fault}
\newacronym{simd}{SIMD}{Single Instruction Multiple Data}
\newacronym{isr}{ISR}{Interrupt Service Routine}
% \newacronym{gpf}{gPF}{Guest Page Fault}
% \newacronym{npt}{NPT}{Nested Page Table}
% \newacronym{seves}{SEV-ES}{SEV-Encrypted State}
% \newacronym{asid}{ASID}{Address Space Identifier}
% \newacronym{sevsnp}{SEV-SNP}{SEV-Secure Nested Paging}
% \newacronym{vmcb}{VMCB}{\gls{vm} Control Block}
% \newacronym{oram}{ORAM}{Oblivious RAM}
% \newacronym{vmsa}{VMSA}{\gls{vm} Save Area}
% \newacronym{rmp}{RMP}{Reverse Mapping Table}
% \newacronym{gva}{gVA}{Guest Virtual Address}
% \newacronym{gpa}{gPA}{Guest Physical Address}
% \newacronym{gpt}{gPT}{Guest Page Table}
\newacronym{hpt}{hPT}{Host Page Table}
% \newacronym{hpa}{hPA}{Host Physical Address}
% \newacronym{tcb}{TCB}{Trusted Computing Base}
% \newacronym{sme}{SME}{Secure Memory Encryption}
% \newacronym{sp}{SP}{Secure Processor}
% \newacronym{vek}{VEK}{VM Encryption Key}
\newacronym{irq}{IRQ}{Interrupt Request}
\newacronym{tsx}{TSX}{Transactional Synchronization Extensions}
\newacronym{smp}{SMP}{Symmetric Multiprocessing}
% \newacronym{slat}{SLAT}{Second-level Address Translation}
\newacronym{ume}{UME}{Universal Memory Engine}
\newacronym{pamt}{PAMT}{Physical-Address-Metadata Table}
\newacronym{uos}{UOS}{Unobservable Storage}
\newacronym{posmap}{PosMap}{Position Map}
\newacronym{pud}{PUD}{Page Upper Directory}
\newacronym{pmd}{PMD}{Page Middle Directory}
\newacronym{pgd}{PGD}{Page Global Directory}
\newacronym{ptept}{PTEPT}{PTE Page Table}
\newacronym{sept}{SEPT}{Secure Extended Page Table}
\newacronym{apic}{APIC}{Advanced Programmable Interrupt Controller}
\newacronym{pit}{PIT}{Programmable Interval Timer}

% \newacronym{vma}{VMA}{Virtual Memory Area}
\newacronym{opfh}{o-PFH}{Oblivious Page Fault Handler}

\newacronym{npfgpt}{NPF-on-GPT-Walk}{Nested Page Faults on Guest Page Tables}

\newacronym{pao}{PAO}{Page Access Obliviousness}


\newacronym{ovma}{\textit{O}-VMA}{Oblivious Virtual Memory Area}
\newcommand*{\apage}{\textit{A}-Page\xspace}
\newacronym{apage}{\apage}{Active Page}
\newacronym{asa}{\textit{A}-Area}{Active subpage area}
\newacronym{pagepool}{\textit{O}-Pool}{ORAM-backed subpage pool}
\newacronym{pager}{\textit{O}-Pager}{Oblivious Pager}

\newacronym{gpfn}{GPFN}{Guest Physical Frame Number}

\newacronym{aex}{AEX}{Asynchronous Enclave Exit}
\newcommand*{\oone}{\textbf{O1}}
\newcommand*{\otwo}{\textbf{O2}}
\newcommand*{\othree}{\textbf{O3}}
\newcommand*{\veriint}{\textsc{VeriInt}}


\newcommand*{\chalfullsys}{\textbf{C1}}
\newcommand*{\chaltimer}{\textbf{C2}}
\newcommand*{\chaloverhead}{\textbf{C3}}

\newcommand*{\optimer}{\textbf{O1}}
\newcommand*{\opmmu}{\textbf{O2}}

\newcommand*{\objbin}{\textbf{OBJ1}}
\newcommand*{\objss}{\textbf{OBJ2}}

\newcommand*{\tickrateavg}{\text{$R_{avg}(Tick)$}}
\newcommand*{\tickratetarget}{\text{$R_{target}(Tick)$}}


\newcommand{\actrg}[1]{$\mathcal{AR}_{#1}$\xspace}
% \newcommand{\actrg}[2]{\ensuremath{AR_{\text{#1}}^{\text{#2}}}\xspace}

% \newcommand{\activeregion}{$\mathcal{A}$\xspace}
% \newcommand{\activecode}{$\mathcal{A}_{code}$\xspace}
% \newcommand{\activedata}{$\mathcal{A}_{data}$\xspace}
% \newcommand{\activept}{$\mathcal{A}_{PT}$\xspace}

\newcommand{\nbench}{\texttt{nbench}\xspace}

\newcommand{\confbaseline}[0]{\textsc{Baseline}\xspace}
\newcommand{\confopaging}[0]{\textsc{+O-Paging}\xspace}
\newcommand{\confsched}[0]{\textsc{+Sched}\xspace}
\newcommand{\confadaptive}[0]{\textsc{+Adaptive}\xspace}
\newcommand{\confnonadaptive}[0]{\textsc{+NonAdaptive}\xspace}

\newcommand{\lnormal}[0]{L1\xspace}
\newcommand{\lnginx}[0]{L2\xspace}
\newcommand{\lredis}[0]{L3\xspace}
\newcommand{\llenpf}[0]{L4\xspace}
\newcommand{\lnpf}[0]{L5\xspace}
\newcommand{\lss}[0]{L6\xspace}


\newcommand{\cntvmexitasync}[0]{\ensuremath{\text{VMExit}_{\text{Async}}}\xspace}


% Frequency macros
\newcommand{\freq}[2]{\ensuremath{f_{\text{#1}}^{\text{#2}}}\xspace}
\newcommand{\freqvmexit}[0]{\ensuremath{f_{\text{VMExit}}}\xspace}
\newcommand{\freqvmexitof}[1]{\ensuremath{f_{\text{VMExit}}^{\text{#1}}}\xspace}
\newcommand{\freqrr}{\ensuremath{f_{\text{ReRand}}}\xspace}
% \newcommand{\freqrrof}[1]{\ensuremath{f_{\textit{\hspace{-0.1em}\scalebox{0.8}{ReRand}}}^{\text{\raisebox{0.2ex}{\scalebox{0.8}{#1}}}}}\xspace}
\newcommand{\freqrrof}[1]{\ensuremath{f_{\text{ReRand}}^{\text{#1}}}\xspace}
\newcommand{\freqrrofof}[2]{\ensuremath{f_{\text{#1}}^{\text{#2}}}\xspace}

\newcommand{\freqsmp}{\ensuremath{f_{\text{Sample}}}\xspace}
\newcommand{\freqsmpof}[1]{\ensuremath{f_{\text{Sample}}^{\text{#1}}}\xspace}

% Period macros
\newcommand{\periodvmexit}{\ensuremath{T_{\text{VMExit}}}\xspace}
\newcommand{\periodvmexitof}[1]{\ensuremath{T_{\text{VMExit}}^{\text{#1}}}\xspace}
\newcommand{\periodrr}{\ensuremath{T_{\text{ReRand}}}\xspace}
\newcommand{\periodrrof}[1]{\ensuremath{T_{\text{ReRand}}^{\text{#1}}}\xspace}

\newcommand{\rerandfunc}[0]{\ensuremath{\mathbb{F}_{\text{Alarm}}}\xspace}

\newcommand{\reqtransparent}{\textbf{R1}\xspace}
\newcommand{\reqsovereign}{\textbf{R2}\xspace}
\newcommand{\reqrobustpolicy}{\textbf{R3}\xspace}
\newcommand{\reqfullsystem}{\textbf{R4}\xspace}
\newcommand{\falarmthreshold}{\textcolor{red}{XXX.XX (Please set me)}}

\newcommand{\AR}[2]{$\mathcal{AR}_{#1}^{\raisebox{-0.20ex}{\scriptsize $\mathnormal{#2}$}}$}


\renewcommand{\thename}{\incognitos}

\renewcommand{\key}[1]{#1}
\section{Abstract}

Recent works have repeatedly proven the practicality of side-channel attacks in undermining the confidentiality guarantees of Trusted Execution Environments such as Intel SGX.
Meanwhile, the trusted execution in the cloud is witnessing a trend shift towards confidential virtual machines (CVMs).
Unfortunately, several side-channel attacks have survived the shift and are feasible even for CVMs, along with the new attacks discovered on the CVM architectures.
Previous works have explored defensive measures for securing userspace enclaves (i.e., Intel SGX) against side-channel attacks.
However, the design space for a CVM-based obfuscation execution engine is largely unexplored.

This paper proposes a unikernel design named \thename to provide full-system obfuscation for CVM-based cloud workloads.
\thename fully embraces unikernel principles such as minimized TCB and direct hardware access to render full-system obfuscation feasible.
\thename retrofits two key OS components, the scheduler and memory management, to implement a novel adaptive obfuscation scheme.
\thename's scheduling is designed to be self-sovereign from the timer interrupts from the untrusted hypervisor with its synchronous tick delivery.
This allows \thename to reliably monitor the frequency of the hypervisor's possession of execution control (i.e., VMExits) and adjust the frequency of memory rerandomization performed by the paging subsystem, which transparently performs memory rerandomization through direct MMU access.
The resulting \thename design makes a case for a self-obfuscating unikernel as a secure CVM deployment strategy while further advancing the obfuscation technique compared to previous works.
Evaluation results demonstrate \thename's resilience against CVM attacks and show that its adaptive obfuscation scheme enables practical performance for real-world programs.

% \pagenumbering{arabic}


